\subsection{用数值密度定义的测度}

定义 2. --- 设 f 是定义在 T 上、取值于 \(\overline{R}\) 或 Banach 空间的函数。若称 f 是局部 \(\mu\) -可积的，意指 f 是 \(\mu\) -可测的，并且 T 中每一点 \(x \in T\) 都有一个邻域 V，使得 \(\mu^{\bullet}(|\mathbf{f}|\varphi_{\mathrm{V}}) < +\infty\) 。

这个定义在 T 为局部紧的情况下与 Ch. V, §5, No.~1 中给出的定义一致。

设 \(f\) 是一个局部 \(\mu\) -可积的正函数；映射 \(K \mapsto f_K \cdot \mu_K\) 是一个预测度（Ch. V, §7, No.~1, Cor. 2 of Th. 1），将其记作 \(f \cdot \mu\) 。

命题 2. --- 若 \(f\) 是正的局部 \(\mu\) -可积函数，则对于每个函数 \(g \in \mathcal{F}_{+}(\mathrm{T})\) ，都有关系式

\[
(f \cdot \mu) ^ {\bullet} (g) = \mu^ {\bullet} (f g).\tag{1}
\]

事实上，对 T 中的每个紧集 K，

\[
(f \cdot \mu) _ {\mathrm{K}} ^ {\bullet} (g _ {\mathrm{K}}) = (f _ {\mathrm{K}} \cdot \mu_ {\mathrm{K}}) ^ {\bullet} (g _ {\mathrm{K}}) = \mu_ {\mathrm{K}} ^ {\bullet} (f _ {\mathrm{K}} g _ {\mathrm{K}}) = \mu_ {\mathrm{K}} ^ {\bullet} ((f g) _ {\mathrm{K}}),
\]

利用 \(f \cdot \mu\) 的定义以及 Ch. V, §5, No.~3 的 Prop. 3。将 \(K\) 取遍并取上确界，即得 Prop. 2。

现在令 \(x \in \mathbf{T}\)，并令 \(\mathbf{V}\) 是 \(x\) 的一个邻域，使得 \(\mu^{\bullet}(f\varphi_{\mathrm{V}}) < +\infty\)（Def. 2）；于是 \((f \cdot \mu)^{\bullet}(\mathrm{V}) = \mu^{\bullet}(f\varphi_{\mathrm{V}}) < +\infty\) ，因此 \(f \cdot \mu\) 是一个测度。

定义 3. --- 设 \(f\) 是一个局部 \(\mu\) -可积的正函数。测度 \(f \cdot \mu : K \mapsto f_K \cdot \mu_K\) 称为具有密度 \(f\) 且相对于 \(\mu\) 的测度，或 \(\mu\) 与函数 f 的乘积测度。每个形如 \(f \cdot \mu\) 的测度（其中 f 为正的局部 \(\mu\) -可积函数）称为以 \(\mu\) 为基的测度。

注记。--- 1) \(f \cdot \mu\) 的定义可以推广到 \(f\) 是复的局部可积函数的情形；此时有 \(|f \cdot \mu| = |f| \cdot \mu\) ，这立即蕴含 \(f \cdot \mu\) 是测度而不只是预测度。我们保留“以 \(\mu'\) 为基的测度”这一表述，用来指称如此定义的复测度。

\begin{enumerate}
\def\labelenumi{\arabic{enumi})}
\setcounter{enumi}{1}
\tightlist
\item
  类似地，若 \(\theta\) 是一个复测度，则称 \(f\) 局部 \(\theta\) -可积，如果 f 局部 \(|\theta|\) -可积；并定义测度 \(f \cdot \theta : K \mapsto f_K \cdot \theta_K\) 。有 \(|f \cdot \theta| = |f| \cdot |\theta|\)（Ch. V, §5, No.~2, Prop. 2）。本号中，我们将撇开一切关于非正测度的内容。
\end{enumerate}

命题 3. --- 设 \(\nu\) 是 T 上的测度。要使 \(\nu\) 形如 \(f \cdot \mu\)，其中 \(f\) 是局部 \(\mu\) -可积的正函数，必要且充分的条件是每个 \(\mu\) -可忽略的紧集都是 \(\nu\) -可忽略的。若 \(f'\) 是另一个局部 \(\mu\) -可积的函数且 \(\nu = f' \cdot \mu\)，则 \(f = f'\) 局部 \(\mu\) -几乎处处成立。

这个条件显然是必要的（Prop. 2）。反之，假设每个 \(\mu\) -可忽略的紧集都是 \(\nu\) -可忽略的。引入 \((\mathrm{K}_{\alpha})_{\alpha \in \mathrm{A}}\) 作为 \(\mathbf{T}\) 关于测度 \(\mu + \nu\) 的一个分割，并令 \(\mathbf{N} = \mathbf{T} - \bigcup_{\alpha \in \mathrm{A}} \mathrm{K}_{\alpha}\) 。显然，\((\mathrm{K}_{\alpha})_{\alpha \in \mathrm{A}}\) 也是 \(\mu\) 和 \(\nu\) 的分割，因此 §1, No.~8 的 Prop. 9 蕴含，对于每个 \(g \in \mathcal{F}_+\) 有以下关系：

\[
\mu^ {\bullet} (g) = \sum_ {\alpha \in A} \mu_ {K _ {\alpha}} ^ {\bullet} \left(g _ {K _ {\alpha}}\right), \quad \nu^ {\bullet} (g) = \sum_ {\alpha \in A} \nu_ {K _ {\alpha}} ^ {\bullet} \left(g _ {K _ {\alpha}}\right).
\]

考虑一个紧集 \(C \subset K_{\alpha}\)，它对于 \(\mu_{K_{\alpha}}\) 是可忽略的；则 C 在局部意义下对于 \(\mu\) 可忽略，从而对于 \(\nu\) 可忽略，最后按照 \(\nu_{K_{\alpha}}\) 的定义，它对于 \(\nu\) 可忽略。由 Lebesgue--Nikodym 定理（Ch. V, §5, No.~5, Th. 2）可知，\(\nu_{K_{\alpha}}\) 关于 \(f_{\alpha}\) 具有一个密度 \(\mu_{K_{\alpha}}\)。令 f 为如下函数：在每个集合 \(f_{\alpha}\) 上与 \(K_{\alpha}\) 一致，在 N 上为 0；则 f 是 \(\mu\) -可测的（Ch. IV, §5, No.~10, Prop. 16），并且对于每个函数 \(g \in F_{+}\)，由上述关系及 Ch. V, §5, No.~3 的 Prop. 3，有

\[
\nu^ {\bullet} (g) = \sum_ {\alpha \in \mathrm{A}} \nu_ {\mathrm{K} _ {\alpha}} ^ {\bullet} \left(g _ {\mathrm{K} _ {\alpha}}\right) = \sum_ {\alpha \in \mathrm{A}} \mu_ {\mathrm{K} _ {\alpha}} ^ {\bullet} \left(f _ {\alpha} g _ {\mathrm{K} _ {\alpha}}\right) = \sum_ {\alpha \in \mathrm{A}} \mu_ {\mathrm{K} _ {\alpha}} ^ {\bullet} ((f g) _ {\mathrm{K} _ {\alpha}}) = \mu^ {\bullet} (f g).
\]

首先可得 \(f\) 局部 \(\mu\) -可积：若 \(x\) 是 \(T\) 中的一点，且 \(V\) 是 \(x\) 的一个邻域并满足 \(\nu^{\bullet}(V) < +\infty\)，则 \(\mu^{\bullet}(f\varphi_V) < +\infty\)。其次，Prop. 2 表明测度 \(\nu\) 与 \(f \cdot \mu\) 具有相同的本质上积分。因此二者相等（§1, No.~2, Cor. of Prop. 2）。\(f\) 的唯一性由紧空间的情形显然可得，命题得证。

注记 3). --- 关于测度 \(\nu = f \cdot \mu\) 的积分理论立即归结为 Ch. V 中处理的理论。为此，令 \((\mathrm{K}_{\alpha})_{\alpha \in \mathbf{A}}\) 是 T 关于 \(\mu\) 的一个分割，因而也是关于 \(\nu\) 的分割，并令 \(\mathrm{T}'\) 是 §1, No.~8 的 Scholium 中定义的局部紧空间；我们把 \(\mu\)（相应地，\(\nu\)）关联到测度 \(\mu'\)（相应地，\(\nu'\)），这些测度定义在 \(\mathrm{T}'\) 上，使得对于 \(\mu\) 和 \(\mu'\)（相应地，对于 \(\nu\) 和 \(\nu'\)），可测函数、取值于 Banach 空间的本质可积函数以及正函数的本质上积分都相同。因此 \(f\) 是 \(\mu'\) -可测的；f 是局部 \(\mu'\) -可积的，因为 \(\mathrm{T}'\) 是局部紧的，并且 \(\mathrm{T}'\) 的一个紧子集只与有限个紧集 \(\mathrm{K}_{\alpha}\) 相交（\(\alpha \in \mathbf{A}\)）。最后，关系 \(\nu'^{\bullet}(g) = \nu^{\bullet}(g) = \mu^{\bullet}(fg) = \mu'^{\bullet}(fg)\) 证明 \(\nu' = f \cdot \mu'\)（Ch. V, §5, No.~3, Prop. 3）。把 Ch. V, §5 的结果转写出来的工作留给读者。

